This appendix contains pseudocode for the main algorithm used to do dataset
selection for DataComp. It also contains additional implementation details on
how metagradients were applied to CLIP, and how they were specifically applied
to the DataComp setting.

\subsection{Dataset Selection Using MGD}

When implementing Algorithm \ref{alg:depsing}, there are several differences from the pseudocode below:
firstly, rather than selecting $\mathbf{m}$ fully randomly every step, we
randomly select a shard comprising fraction $p$ of the data and take steps on
all datapoints in the shard (see Section~\ref{sec:scaling_clip}).  To mitigate
overfitting, we also bake a ``minibatch fraction'' $q$ into our model output
function $\phi$.  For example, if $\phi$ calculates model loss on the ImageNet
train set, each time $\phi$ is called, we randomly sample fraction $q$ of the
ImageNet train set to evaluate on.

\paragraph{Adapting the CLIP loss function to our surrogate learning algorithm.}
Here, we explain how dataweights are incorporated into the CLIP loss
function---the formulation given in Section \ref{sec:datacomp} is actually
slightly simplified and incorrect, as it does not account for cross terms in
the CLIP contrastive loss.  As a refresher, we first state the ``vanilla'' CLIP
loss function, $\ell$, as it is defined in \cite{radford2021learning}.  Letting
$b$ be the batch size and $d$ be the embedding dimension, and $\mathbf{x}$ be
the
training batch at timestep $k$.
Recall that the CLIP model internally has two ``submodules'': and image
embedder, and a text embedder.  We then use these to obtain image embeddings
$E_I \in \mathbb{R}^{b \times d}$ and text embeddings $E_T \in
\mathbb{R}^{b \times d}$ from $\mathbf{x}$. We then compute the image-wise
scores, or logits, for this batch as
$S = E_I E_T^\top$ \footnote{The CLIP model scales these logits by a
	temperature parameter $\tau$ before applying the softmax. While we omit
	$\tau$ in our definitions, it can be easily incorporated. All our
experiments use temperature scaling.}.
Then, we can define the CLIP loss (as a function of the logits) as
\begin{equation*}
	L(S) = \frac{1}{2} (L_I(S) + L_T(S)),
\end{equation*}
where $L_I$ and $L_T$ are row-wise and column-wise cross-entropy losses, respectively:
\begin{equation*}
	L_I(S) = \sum_{i=1}^b 
    \log \left(
      \frac{\exp(S_{i,i})}
           {\sum_{j=1}^b \exp(S_{i,j})}
    \right), \quad
	L_T(S) = \sum_{i=1}^b 
    \log \left(
      \frac{\exp(S_{i,i})}
           {\sum_{j=1}^b \exp(S_{j,i})}
    \right).
\end{equation*}

We now wish to relax $L$ into a new function $L'$ that supports an additional
input $\mathbf{z} \in \mathbb{R}^n$, where $\frac{\partial L'}{\partial
\mathbf{z}}$ resembles the metagradients with respect to dataweights.  In order
to do this, we imagine expanding passing the \emph{entire dataset} $D$ into our
embedder to obtain $E_I'$ and $E_T'$, and take our new logits $S' = E_I'
E_T'^\top \in \mathbb{R}^{n \times n}$.

There are some additional key conditions our relaxation $L'$ should satisfy. 
Particularly:  when $\mathbf{z} = \mathbf{0}_n$, we should recover the normal
CLIP loss $L$, and when $\mathbf{z}$ is all 0's except for a single entry $i$,
$L'$ should act as if $i$ had been appended to the original batch $\mathbf{x}$.
In addition, $L'$ should always have meaningful partials with
respect to $\mathbf{z}$, even when some values in $\mathbf{z}$ are 0.  

Letting $\mathbf{1}_{i = j}$ and $\mathbf{1}_{i \neq j}$ be indicator variables
and letting $\mathbf{1}_{k} \in \{0, 1\}^n$ be the indicator vector for the
$k$-th batch, we find that the definition
\[
	L'(S', \mathbf{z}) = L'_I(S', \mathbf{z}) + L'_T(S', \mathbf{z}),
\]
where
\[
	L'_I(S', \mathbf{z}) = \sum_{i=1}^n
	(z_i + (\mathbf{1}_{k})_i) \log \left(
      \frac{\exp(S'_{i,i})}
	  {\sum_{j=1}^n \exp(S'_{i,j}) \left(\mathbf{1}_{i = j} + \mathbf{1}_{i \neq j} (z_j + (\mathbf{1}_k)_j)\right)}
    \right)
\]
and
\[
	L'_T(S', \mathbf{z}) = \sum_{i=1}^b 
	(z_i + (\mathbf{1}_{k})_i) \log \left(
      \frac{\exp(S'_{i,i})}
	  {\sum_{j=1}^n \exp(S'_{j,i}) \left(\mathbf{1}_{i = j} + \mathbf{1}_{i \neq j} (z_j + (\mathbf{1}_k)_j)\right)}
    \right)
\]
satisfy these conditions.  

Finally, we let define the loss for the entire batch $\ell'$ as a function of $\mathbf{z}$ and model parameters $\theta$ which outputs the loss calculated according to $L'$ above.
To summarize, letting $\mathbf{x}^{(t)}$ denote the $t$-th training batch, the loss function $\ell_t$ at step $t$ of our surrogate learning algorithm $\mathcal{A}'$ for CLIP training is:
\begin{equation*}
	\label{eq:surrogate_objective_new}
    \ell'_t(\theta) := 
    \begin{cases}
		\ell(\mathbf{x}^{(t)}; \theta) & \text{if } t \neq k \\
		\ell'(\mathbf{z}; \theta) & \text{if } t = k.
    \end{cases}
\end{equation*}
We find that this empirically works well for obtaining meaningful
metagradients with respect to dataweights in the CLIP setting, and yields to strong dataset selection results.

\subsection{Scaling MGD for CLIP and DataComp}
\label{sec:scaling_clip}

MGD is highly scalable, allowing it to be applied to large-scale settings like
training CLIP models.  In particular, computing metagradients is only up to a
constant factor more expensive than training a model normally. 
Here, we outline challenges we faced in scaling MGD in this setting, and how they were resolved.
Specifically, we will explain how we efficiently calculated metagradients for
CLIP models and efficiently tracked/shuffled our dataset selection from
step-to-step despite its large storage footprint.

\paragraph{Computing metagradients.}
Due to the large batch size used in the CLIP contrastive loss, we implement
manual gradient checkpointing to make the operations computationally feasible
on our hardware.  The most memory-intensive operation are model forward passes
(and its gradients): obtaining the image and label embeddings given raw pixel
data and tokens.  So, we manually make gradient checkpoints before this
operation, allowing us to run the embedder in minibatches to avoid memory
issues.  This setup also naturally lends itself to parallelization across
multiple GPU's, which we make use of to further speed up our computations.

\paragraph{Loading, writing, and storing data.}
Due to the data-intensive nature of training large models like CLIP and our
need to frequently produce new datasets at each optimization step, we found
that using the webdataset \cite{webdataset2024webdataset} format given by
DataComp was restrictively slow.  To circumvent this, we rewrote all data
following the format of FFCV \cite{leclerc2022ffcv}, allowing us to load and
write data much faster.  Specifically, we divided the entire candidate pool
into 8 base shards.  Once we trained a model, we choose one of the 8 shards,
compute metagradients corresponding to all datapoints in the shard, take a
gradient step on them, and rewrite the shard.  This roughly corresponds to $p =
\frac{1}{8}$ in Algorithm \ref{alg:depsing}, which we empirically worked well
for optimizing.  In following steps, we always choose one of the 8
\emph{original} shards to calculate metagradients for---this ensures that
points removed from the dataset in some optimization step can return if they
have a negative metagradient.

We also observed that always stepping on the sign causes the sizes of the
shards to grow over time: stepping based on the sign of the metagradient does
not decrease the weight on a positive-weight datapoint if its dataweight  is
already 0, so our steps are biased towards increasing the size of the shards. 
To combat this blowup, after some number of optimization steps, we choose a
fixed shard size and enforce that subsequent steps must not change the size of
the shards---the step size thereafter is controlled by hyperparameter $q$
representing the fraction of datapoints in a shard which are incremented.  We
experimented both with randomly sampling which points are added or removed, and
stepping on the datapoints with the top $q$ and bottom $q$ metagradients; the 
latter seems to give empirically better performance.

To maintain randomness during shuffling, we implement an 8-way dataloader which
would shuffle all 8 shards individually.  Then, to sample a batch of $b$
datapoints, we would sample $b / 8$ datapoints from each shard and concatenate
them to fill our batch.  This works better than simply sampling our entire
batch from a single shard, as (especially in later optimization steps) shards
may contain a high number of duplicate datapoints, which causes CLIP's
contrastive loss function to misbehave if they appear in the same batch.

To minimize disk space used, old shards can be deleted once they become
``stale''.  Specifically, if shard $s$ is rewritten into shard $s'$, all future
optimization steps will never read $s$ again, and $s$ can safely be
deleted.  Thus, when running MGD for a large number of steps and potentially
rewriting each shard multiple times, the total disk space used by our algorithm
is constant in the number of steps we take: it stores the 8 most recently written 
shards on disk at any given time, and any other shards are deleted to save space.

\subsection{Details Pertaining to the DataComp Benchmark}

\paragraph{Setting.}
We provide a brief summary of the DataComp competition here, and we refer
readers to the original paper \cite{gadre2024datacomp}.  DataComp is a
framework to compare different training dataset selection techniques. 
Participants submit a training dataset (which, for our purposes, is a subset of
a larger dataset), upon which a CLIP model is trained from scratch with a fixed
learning algorithm, model architecture, and number of training steps.  We focus
on DataComp-small, which has a candidate pool of 12.8 million samples.  The
number of training steps in this case is also fixed at 12.8 million samples. 

We try to match the optimization hyperparameters enforced by DataComp as
closely as possible.  As a refresher, our ADAM\cite{kingma2015adam} update step
can be written as
\begin{equation}
	\label{eq:adam_update}
	\theta_{t+1} = -\alpha_t \cdot \left(m_t / \left(\sqrt{v_t +
	\epsilon_\mathrm{root}} + \epsilon\right) + \lambda \theta_t\right)
\end{equation}
where $m_t$ and $v_t$ are running estimates of the first and second moments of
the gradients, respectively, $\lambda$ represents weight decay, $\alpha$
represents the learning rate, and $\epsilon$ and $\epsilon_\mathrm{root}$ are
hyperparameters to avoid blowup.  Our training hyperparameters can be found in
Table \ref{tab:clip_hparams} and are identical to those mandated by
DataComp-small, aside from a positive $\epsilon_\mathrm{root}$ added for
numerical stability. The values of $\epsilon_{\mathrm{root}}$ and $k$ (the step
at which metagradients are calculated) were chosen to empirically maximize
metasmoothness.


\begin{table}[h]
	\caption{Hyperparameters for the CLIP DataComp experiments.}
	\centering
	\begin{tabular}{lc}
\toprule
\textbf{Hyperparameter} & \textbf{Value} \\
\midrule
DataComp Scale & small \\
Model & ViT-B/32 \\
Train compute (MACs) & $9.5 \times 10^{16}$ \\
Pool size & 12.8M \\
\# samples seen & 12.8M \\
Batch size & 4096 \\
Training batches & 3125 \\
$k$ & 2800 \\
Learning rate & $5 \times 10^{-4}$ \\
AdamW $\beta_1$ & 0.9 \\
AdamW $\beta_2$ & 0.98 \\
AdamW $\epsilon_{\textrm{root}}$ & $1 \times 10^{-17}$ \\
Warmup & 500 \\
\bottomrule
\end{tabular}

	\label{tab:clip_hparams}
\end{table}

Our experiments are also run on an incomplete subset of the entire DataComp
candidate pool.    DataComp did not store the raw image and text files when
assembling their dataset; they only stored a list of URL's to download data
from.  Due to the nature of the internet, for various reasons, some of these
URL's no longer point to the same data (or no longer point to any data at all).
Thus, after ignoring these broken links, our candidate pool is only around
$80\%$ of the size of the original DataComp candidate pool when it was
collected in 2023.  All our results are obtained by running our methods on this
subset of the DataComp pool.

\paragraph{Evaluation tasks.}
In order to ensure that our method is truly improving trained models'
performances on the \emph{entire target distribution} and not overfitting to
the target set, for each of the 38 evaluation tasks used by DataComp, we
attempted to separately create a disjoint target and validation set (DataComp only
creates test sets for each task).  Thus, metagradients were computed on the
target sets and model performance was evaluated on the validation set, before
submitting with the official DataComp script and evaluating on the test sets. 
This ensures that our method's generalization ability is being evaluated, and
we are not overfitting to our target set.

For various reasons, creating target splits was not possible for all 38 tasks;
we summarize our setup in Table \ref{tab:datacomp_tasks}.

{
\renewcommand{\thefootnote}{\fnsymbol{footnote}}
\begin{table}[h]
\caption{All DataComp evaluation tasks.  The ``Target set'' column refers to whether metagradients were taken on the target set corresponding to this dataset.}
\centering
{
	\scriptsize
	\begin{tabular}{llllllc}
\toprule
\textbf{Dataset}  &  \textbf{Task}  &  \textbf{Test size}  &  \textbf{Train size}  &  \textbf{Val size}  &  \textbf{Main metric}  &  \textbf{Target set} \\
\midrule
Caltech-101 \cite{fei2004learning}  &  Object recognition  & 6085 &  2754  &  306  &  mean per class  &  \checkmark \\
CIFAR-10 \cite{krizhevsky2009learning}  &  Visual recognition  & 10000 &  45000  &  5000  &  accuracy  &  \checkmark \\
CIFAR-100 \cite{krizhevsky2009learning}  &  Visual recognition  & 10000 &  45000  &  5000  &  accuracy  &  \checkmark \\
CLEVR Counts \cite{johnson2017clevr,zhai2019large}  &  Counting  & 15000 &  65000  &  5000  &  accuracy  &  \checkmark \\
CLEVR Distance \cite{johnson2017clevr,zhai2019large}  &  Distance prediction  & 15000 &  65000  &  5000  &  accuracy  &  \checkmark \\
Country211 \cite{radford2021learning,thomee2016yfcc100m}  &  Geolocation  & 21100 &  37980  &  4220  &  accuracy  &  \checkmark \\
DTD \cite{cimpoi2014describing}  &  Texture classification  & 1880 &  3384  &  376  &  accuracy  &  \checkmark \\
EuroSAT \cite{helber2019eurosat,zhai2019large}  &  Satellite imagery recognition  & 5400 &  19440  &  2160  &  accuracy  &  \checkmark \\
FGVC Aircraft \cite{maji2013fine}  &  Aircraft recognition  & 3333 &  6001  &  666  &  mean per class  &  \checkmark \\
Food-101 \cite{bossard2014food}  &  Food recognition  & 25250 &  70750  &  5000  &  accuracy  &  \checkmark \\
GTSRB \cite{stallkamp2011german}  &  Traffic sign recognition  & 12630 &  35289  &  3920  &  accuracy  &  \checkmark \\
ImageNet 1k \cite{deng2009imagenet}  &  Visual recognition  & 50000 &  1276167  &  5000  &  accuracy  &  \checkmark \\
ImageNet Sketch \cite{wang2019learning}  &  Visual recognition  & 50889 &  N/A  &  N/A  &  accuracy  &   \footnotemark[1] \\
ImageNet V2 \cite{recht2019imagenet}  &  Visual recognition  & 10000 &  N/A  &  N/A  &  accuracy  &   \footnotemark[1] \\
ImageNet-A \cite{hendrycks2019natural}  &  Visual recognition  & 7500 &  N/A  &  N/A  &  accuracy  &   \footnotemark[1] \\
ImageNet-O \cite{hendrycks2019natural}  &  Visual recognition  & 2000 &  N/A  &  N/A  &  accuracy  &   \footnotemark[1] \\
ImageNet-R \cite{hendrycks2020faces}  &  Visual recognition  & 30000 &  N/A  &  N/A  &  accuracy  &   \footnotemark[1] \\
KITTI distance \cite{geiger2012ready,zhai2019large}  &  Distance prediction  & 711 &  N/A  &  N/A  &  accuracy  &   \footnotemark[2] \\
MNIST \cite{lecun1998mnist}  &  Digit recognition  & 10000 &  55000  &  5000  &  accuracy  &  \checkmark \\
ObjectNet \cite{barbu2019objectnet}  &  Visual recognition  & 18574 &  N/A  &  N/A  &  accuracy  &  \footnotemark[1] \\
Oxford Flowers-102 \cite{nilsback2008automated}  &  Flower recognition  & 6149 &  1836  &  204  &  mean per class  &  \checkmark \\
Oxford-IIIT Pet \cite{parkhi2012cats,zhai2019large}  &  Pet classification  & 3669 &  3312  &  368  &  mean per class  &  \checkmark \\
Pascal VOC 2007 \cite{everingham2010pascal}  &  Object recognition  & 14976 &  14096  &  1566  &  accuracy  &  \checkmark \\
PatchCamelyon \cite{veeling2018rotation,zhai2019large}  &  Metastatic tissue cls.  & 32768 &  289912  &  5000  &  accuracy  &  \checkmark \\
Rendered SST2 \cite{zhai2019large}  &  Sentiment classification  & 1821 &  7013  &  779  &  accuracy  &  \checkmark \\
RESISC45 \cite{cheng2017remote,zhai2019large}  &  Satellite imagery recognition  & 6300 &  22680  &  2520  &  accuracy  &  \checkmark \\
Stanford Cars \cite{krause20133d}  &  Vehicle recognition  & 8041 &  7329  &  814  &  accuracy  &  \checkmark \\
STL-10 \cite{coates2011analysis}  &  Visual recognition  & 8000 &  4500  &  500  &  accuracy  &  \checkmark \\
SUN-397 \cite{xiao2010sun}  &  Scene recognition  & 108753 &  N/A  &  N/A  &  accuracy  &  \footnotemark[3] \\
SVHN \cite{netzer2011reading,zhai2019large}  &  Digit recognition  & 26032 &  68257  &  5000  &  accuracy  &  \checkmark \\
iWildCam \cite{beery2021iwildcam,koh2020wilds}  &  Animal recognition  & 42791 &  147084  &  5000  &  macro F1 score  &  \checkmark \\
Camelyon17 \cite{bandi2018from,koh2020wilds}  &  Metastatic tissue cls.  & 85054 &  365900  &  5000  &  accuracy  &  \checkmark \\
FMoW \cite{christie2018functional,koh2020wilds}  &  Satellite imagery recognition  & 22108 &  103261  &  5000  &  worst-region acc.  &  \checkmark \\
Dollar Street \cite{rojas2022dollar}  &  Object recognition  & 3503 &  13842  &  1537  &  worst-income top-5 acc.  &  \checkmark \\
GeoDE \cite{ramaswamy2024geode}  &  Object recognition  & 12438 &  44488  &  4943  &  worst-region acc.  &  \checkmark \\
Flickr30k \cite{young2014from}  &  Image and text retrieval  & 31014 &  N/A  &  N/A  &  R@1  &  \footnotemark[4] \\
MSCOCO \cite{lin2014microsoft}  &  Image and text retrieval  & 5000 &  N/A  &  N/A  &  R@1  &  \footnotemark[4] \\
WinoGAViL \cite{bitton2022winogavil}  &  Commonsense association  & 3563 &  N/A  &  N/A  &  Jaccard score  &  \footnotemark[4] \\
\bottomrule
\end{tabular}

}
\label{tab:datacomp_tasks}
\end{table}
\footnotetext[1]{No train or val set exists for this dataset, so we were unable to create disjoint target and val sets.}
\footnotetext[2]{We were unable to use this dataset due to technical difficulties.}
\footnotetext[3]{Both the train and val sets were used by DataComp to make their test set, so we were unable to create disjoint target and val sets.}
\footnotetext[4]{Retrieval tasks were not used for metagradients.}
}


